% Explicit Euler/gradient-descent iteration: x_{k+1}=x_k-0.1 ∇f(x_k).
% f(x,y)=(x-1)^2+4(y+0.5)^2. Coordinates checked analytically.
\clearpage
\section*{Visual laboratory: descend a known objective}
\noindent Consider the strictly convex quadratic
\(f(x,y)=(x-1)^2+4(y+0.5)^2\), with minimizer
\((1,-0.5)\). Gradient descent at step size \(\alpha=0.1\) gives
\(x_{k+1}=x_k-0.2(x_k-1)\) and
\(y_{k+1}=y_k-0.8(y_k+0.5)\).
The trajectory starts at \((-1,1)\), and each plotted height is evaluated
with the same mathematical objective.
\begin{figure}[!htbp]
\centering
\begin{tikzpicture}
\begin{axis}[
 booktwoaxis,width=0.86\textwidth,height=79mm,
 title={Analytic quadratic objective with gradient-descent iterates},
 xlabel={$x$},ylabel={$y$},zlabel={$f(x,y)$},
 xmin=-1.25,xmax=2.2,ymin=-1.3,ymax=1.25,zmin=0,zmax=16.0,
 view={49}{32},colormap name=booktwothermal,
 domain=-1.2:2.2,y domain=-1.3:1.2,samples=30,samples y=27
]
\addplot3[surf,shader=interp,draw=none,opacity=0.83]
{(x-1)^2+4*(y+0.5)^2};
\addplot3[booktwonavy,very thick,mark=*,mark size=1.3pt] coordinates {
 (-1,1,13)
 (-0.6,-0.2,2.92)
 (-0.28,-0.44,1.6528)
 (-0.024,-0.488,1.049152)
 (0.1808,-0.4976,0.67111168)
};
\addplot3[only marks,mark=star,mark size=2.4pt,booktwogold]
coordinates {(1,-0.5,0)};
\end{axis}
\end{tikzpicture}
\caption{Analytically specified optimization landscape and the first five
gradient-descent positions, with \(\alpha=0.1\). The star indicates the exact
minimizer. This is not an optimizer benchmark: it is a deterministic
derivation check.}
\label{fig:book02-optimization-landscape}
\end{figure}
\noindent\textbf{Reconstruct it:} implement the iteration, log objective
and gradient norms, and verify every plotted position. Repeat with
\(\alpha=0.2\), where the \(y\)-component alternates and decays, and
with \(\alpha=0.3\), where it alternates but diverges because
\(\lvert 1-8\alpha\rvert>1\). The stability interval in this direction
is \(0<\alpha<\tfrac14\).
